% KHÔNG \input file này: diễn giải HTTP/RFC gần tài liệu chuẩn, dễ trùng Turnitin.

\section{Nền tảng giao thức web và trao đổi dữ liệu}

Các mục trước tập trung NLP và truy hồi. Hệ thống đồ án còn là ứng dụng web:
trình duyệt nói chuyện với FastAPI bằng HTTP, dữ liệu đi JSON hoặc SSE, xác thực
đi header. Mục này giải thích các khái niệm đó ở mức đủ để đọc Chương~3 mà không
cần tài liệu ngoài, rồi gắn từng khái niệm với chỗ dùng trong sản phẩm.

\subsection{HTTP}

HTTP (Hypertext Transfer Protocol) là giao thức tầng ứng dụng: client gửi yêu
cầu, server trả đáp. Một yêu cầu gồm phương thức (GET, POST, PUT, PATCH, DELETE,
OPTIONS), đường dẫn, phiên bản, header và (tuỳ phương thức) thân. Đáp gồm mã
trạng thái, header và thân. HTTP không giữ trạng thái giữa hai request: mỗi lần
gọi độc lập, trừ khi ứng dụng tự mang ngữ cảnh (cookie, token)~\cite{rfc9110}.

Đồ án không tự cài stack TCP. FastAPI và trình duyệt đã nói HTTP. Việc của đồ án
là chọn đúng phương thức và mã:

\begin{itemize}
    \item GET đọc tài nguyên (danh mục luật, cây Điều, hội thoại). Lặp GET không
    được làm đổi dữ liệu.
    \item POST tạo (đăng ký, tin nhắn, tải tệp, sinh lịch).
    \item PATCH sửa một phần (đổi tên hội thoại, đánh dấu task).
    \item DELETE xóa hội thoại hoặc tài liệu của đúng chủ sở hữu.
    \item OPTIONS xuất hiện khi trình duyệt hỏi CORS trước POST khác origin.
\end{itemize}

HTTPS là HTTP đi trên TLS: mã hoá đường truyền, xác thực máy chủ bằng chứng thư.
Compose nội bộ dùng HTTP thuần vì các container nằm trên một máy. Khi đưa hệ
thống ra Internet, TLS phải đặt trước nginx; JWT trên HTTP thường dễ bị nghe lén.

Phiên bản HTTP/1.1 giữ kết nối (keep-alive) để nhiều request dùng một TCP.
SSE của đồ án dựa đúng cơ chế này: một GET/POST dài, server lần lượt ghi từng
dòng \texttt{data:} mà không đóng socket.

\subsection{Mã trạng thái HTTP dùng trong đồ án}

Mã chia nhóm: 2xx thành công, 4xx lỗi phía client, 5xx lỗi phía server. Đồ án
không dùng hết bảng mã RFC, chỉ một tập có nghĩa nghiệp vụ:

\begin{itemize}
    \item 200: đọc thành công.
    \item 201: tạo tài nguyên (ví dụ tài liệu sau khi tải).
    \item 202: đã nhận việc chạy nền (soát xét hợp đồng); client phải poll.
    \item 204: thành công không thân (ít dùng).
    \item 400: dữ liệu hợp lệ schema nhưng sai nghiệp vụ.
    \item 401: thiếu hoặc hỏng token.
    \item 403: đã xác thực nhưng không đủ quyền admin.
    \item 404: không có tài nguyên \emph{hoặc} không được xem tài nguyên người
    khác (cố ý không trả 403 để khỏi dò tồn tại hội thoại).
    \item 409: xung đột, ví dụ mã số thuế trùng.
    \item 422: Pydantic từ chối schema.
    \item 429: vượt slowapi.
    \item 503: hỏi đáp khi chưa có chỉ mục vector hoặc thiếu khóa mô hình.
\end{itemize}

Việc gắn mã này không phải hình thức. Client TanStack Query phân nhánh theo mã:
401 thì làm mới token; 503 thì hiện chữ ``chưa sẵn sàng'' chứ không vẽ stack
trace.

\subsection{JSON}

JSON (JavaScript Object Notation) là định dạng văn bản: đối tượng
\texttt{\{khoá: giá trị\}}, mảng, chuỗi, số, boolean, null~\cite{rfc8259}. JSON
không có kiểu ngày tháng riêng; đồ án đưa thời điểm dạng ISO-8601 chuỗi. JSON
không có comment. FastAPI và Pydantic xuất/nhập JSON mặc định.

Ưu điểm so với XML trong đồ án: gọn, khớp sẵn JavaScript/TypeScript, đủ cho schema
hội thoại và finding hợp đồng. Nhược: không có namespace; số lớn phải cẩn thận vì
IEEE-754 phía JS. Mã số thuế và số Điều được giữ kiểu chuỗi khi cần tránh mất số 0
đầu.

Mọi API REST của đồ án trả JSON, trừ luồng SSE (text/event-stream) và tải tệp
(multipart). Lỗi cũng JSON: trường \texttt{detail}.

\subsection{CORS}

Trình duyệt áp chính sách cùng nguồn (same-origin): một trang
\texttt{http://localhost:5173} không được đọc đáp từ
\texttt{http://localhost:8023} trừ khi server đích cho phép qua CORS
(Cross-Origin Resource Sharing). Cơ chế: trình duyệt gửi OPTIONS (preflight) kèm
\texttt{Origin}, \texttt{Access-Control-Request-Method}; server trả
\texttt{Access-Control-Allow-Origin} và các header cho phép.

Đồ án khoá CORS theo danh sách origin cấu hình, không dùng \texttt{*}, vì
\texttt{*} không đi cùng \texttt{Authorization} một cách an toàn. Trên Docker,
nginx phục vụ SPA và proxy \texttt{/api} cùng origin nên preflight gần như không
chạy. CORS chủ yếu phục vụ máy phát triển: Vite cổng 5173, API cổng 8023.

\subsection{Unicode, UTF-8 và tiếng Việt}

Unicode gán mỗi ký tự một mã điểm. UTF-8 mã hoá mã điểm thành một đến bốn byte,
tương thích ASCII. Tiếng Việt (dấu thanh, đ, ơ) nằm trong dải Latin mở rộng.
Mọi lớp đồ án (Postgres, JSON, React, tệp corpus) dùng UTF-8. Sai encoding làm
vỡ dấu trên Điều luật --- lỗi người dùng phát hiện ngay.

Tìm kiếm không dấu không phải ``bỏ Unicode''. Đồ án giữ bản UTF-8 đầy đủ để hiển
thị, đồng thời tạo bản đã \texttt{unaccent} chỉ để lập chỉ mục. Hai bản phải cùng
một Điều; nếu lệch, highlight trên UI sẽ trỏ sai vị trí ký tự.

\subsection{Biểu thức chính quy}

Biểu thức chính quy (regular expression) mô tả mẫu chuỗi. Đồ án dùng regex ở
vài chỗ hẹp: fallback tách từ khi underthesea không nạp; nhận số Điều trong câu
hỏi; kiểm tra định dạng mã số thuế ở mức hình thức. Regex không thay parser
luật. Văn bản quy phạm có cấu trúc Chương--Điều, không phải bài toán ``tìm mọi
chuỗi giống luật'' bằng một mẫu.

Lạm dụng regex trên HTML hoặc trên toàn bộ PDF thường mong manh. Vì vậy trích
hợp đồng dùng thư viện PDF/DOCX, không dùng một regex ``lấy điều khoản''.

\section{Luồng thực thi máy chủ: đồng bộ, bất đồng bộ, ASGI}

\subsection{I/O chờ mạng}

Một yêu cầu hỏi đáp phần lớn thời gian là chờ: nhúng vector, gọi Gemini, đọc
Postgres, đọc Chroma. Nếu mỗi request chiếm một luồng hệ điều hành và ngủ khi
chờ, số kết nối đồng thời bị giới hạn nhanh. Mô hình bất đồng bộ
(\texttt{async}/\texttt{await}) nhả vòng sự kiện khi chờ I/O, phục vụ request
khác trên cùng tiến trình.

Đồ án: router FastAPI và SQLAlchemy async. Riêng truy hồi lai (BM25 + Chroma) là
mã đồng bộ nặng CPU/I/O đĩa, nên được đẩy sang luồng phụ
(\texttt{asyncio.to\_thread}) để heartbeat SSE không bị nghẽn.

\subsection{WSGI và ASGI}

WSGI là giao diện đồng bộ cổ điển giữa máy chủ Python và ứng dụng (Flask, Django
truyền thống). ASGI bổ sung sự kiện bất đồng bộ và WebSocket/SSE. FastAPI chạy
trên ASGI (Uvicorn/Starlette)~\cite{fastapi}. Nếu gắn FastAPI vào máy chủ WSGI
thuần, SSE và \texttt{async} mất ý nghĩa.

\subsection{OpenAPI}

OpenAPI mô tả API bằng JSON/YAML: đường, phương thức, schema, mã lỗi. FastAPI
sinh OpenAPI từ chữ ký hàm và Pydantic; giao diện Swagger tại \texttt{/docs} dùng
khi trình bày và khi frontend đối chiếu tên trường. Đồ án không sinh client từ
OpenAPI tự động (chưa cần codegen); TypeScript viết tay bám schema.

\section{Tầng trình duyệt: Node.js, npm, CSS, SPA}

\subsection{JavaScript, ECMAScript và môi trường thực thi}

JavaScript là ngôn ngữ của trình duyệt. ECMAScript là chuẩn ngôn ngữ. Node.js
chạy JavaScript phía máy phát triển: cài gói, chạy Vite, chạy Vitest. Sản phẩm
Docker frontend \emph{không} chạy Node khi phục vụ người dùng: chỉ nginx và tệp
tĩnh đã build. Phân biệt này tránh nhầm ``React cần Node trên server''.

\subsection{npm}

npm là trình quản lý gói hệ sinh thái Node. \texttt{package.json} khai báo phụ
thuộc; \texttt{package-lock.json} khóa phiên bản. Đồ án cài React, Vite, Tailwind,
TanStack Query, Zustand, React Router, Axios, Vitest qua npm. Không cam kết thư
mục \texttt{node\_modules} vào báo cáo hay image thừa: image frontend đa giai
đoạn, giai đoạn cuối chỉ còn bản build.

\subsection{Mô hình hộp CSS}

Mọi phần tử HTML là một hộp: nội dung, padding, border, margin. Flexbox và Grid
sắp xếp hộp theo hàng/cột. Tailwind không thay mô hình hộp; Tailwind chỉ là cách
gắn utility class thay vì viết bộ chọn CSS dài. Đồ án chọn Tailwind để giữ giao
diện trong JSX, giảm file CSS toàn cục, phù hợp một người làm cả frontend.

\subsection{Virtual DOM và reconciliation}

React giữ cây mô tả giao diện (element) và đối chiếu với lần render trước để
vá DOM thật tối thiểu~\cite{react}. Khi token SSE tới, đồ án cộng chuỗi vào
state tin nhắn; React re-render khung chat. Nếu ghi thẳng vào
\texttt{innerHTML} từng token, dễ lệch trạng thái và mất kiểm soát XSS.

\subsection{Bundle, tree-shaking và code splitting}

Vite/Rollup gom module thành tệp nhỏ hơn, loại mã không được import
(tree-shaking). Route quản trị \texttt{lazy()} tách bundle: người không vào
admin không tải mã lab. SPA vẫn phải tải runtime React lần đầu; chấp nhận được
vì đối tượng là nhân sự/kế toán sau đăng nhập, không phải trang marketing.

\section{Cơ sở dữ liệu quan hệ: giao dịch, chỉ mục, nhóm kết nối}

\subsection{Mô hình quan hệ và toàn vẹn}

CSDL quan hệ tổ chức dữ liệu thành bảng, cột, khóa chính, khóa ngoại. Toàn vẹn
tham chiếu ngăn hội thoại trỏ user không tồn tại. Đồ án để PostgreSQL bắt các
ràng buộc này, không chỉ tin mã Python.

\subsection{ACID}

Một giao dịch quan hệ lý tưởng có bốn tính chất:

\begin{itemize}
    \item Atomicity: hoặc tất cả lệnh trong giao dịch thành công, hoặc không
    lệnh nào có hiệu lực.
    \item Consistency: ràng buộc schema không bị phá.
    \item Isolation: giao dịch song song không đọc trạng thái dở.
    \item Durability: đã commit thì còn sau khi mất điện (trong giới hạn cấu
    hình WAL của Postgres).
\end{itemize}

Đồ án dùng giao dịch khi đăng ký kèm tạo tổ chức và khi sinh lịch: không để user
tồn tại mà tổ chức dở. RAG đọc corpus không cần giao dịch dài; đọc Điều là truy
vấn ngắn.

\subsection{Chỉ mục B-tree và GIN}

Chỉ mục giảm số hàng phải quét. B-tree phù hợp so sánh bằng và khoảng trên khóa
(id, user\_id, created\_at). GIN phù hợp \texttt{tsvector} và \texttt{pg\_trgm}:
một hàng liên kết nhiều token. Đồ án: B-tree trên khóa ngoại và trường lọc; GIN
trên cột tìm kiếm. Thiếu chỉ mục làm cây văn bản và FTS chậm khi corpus
\soDieu{} Điều; thừa chỉ mục làm ghi chậm hơn --- chấp nhận vì corpus gần như
chỉ đọc sau khi nạp.

\subsection{Nhóm kết nối (connection pooling)}

Mở TCP tới Postgres đắt. Pool giữ sẵn vài kết nối, cho mượn khi có request, trả
lại khi xong. SQLAlchemy async dùng pool. Số kết nối tối đa phải nhỏ hơn
\texttt{max\_connections} của Postgres, đặc biệt khi chạy pytest song song và
Compose cùng lúc. Sai cấu hình pool biểu hiện là timeout lúc demo, không phải
lỗi RAG.

\subsection{Quyền trong CSDL}

User ứng dụng không được \texttt{CREATE EXTENSION}. Extension
\texttt{unaccent}, \texttt{pg\_trgm} nằm \texttt{init.sql} chạy bằng user quản
trị lúc khởi tạo volume. Tách quyền này tránh script backend đòi superuser trên
môi trường thật.

\section{Bảo mật ứng dụng: băm, token, lỗ hổng trình duyệt}

\subsection{Băm một chiều và mã hoá hai chiều}

Mã hoá (encryption) có chìa để giải ra bản rõ --- dùng khi cần đọc lại (ví dụ
sao lưu). Băm (hash) một chiều: từ mật khẩu ra chuỗi cố định, không có chìa giải
ngược. Mật khẩu phải băm, không mã hoá. SHA-256 quá nhanh nên bị dò từ điển bằng
GPU. Argon2 cố ý chậm và tốn bộ nhớ~\cite{argon2}. Đồ án: mật khẩu qua argon2id;
JWT ký HS256 (đây là mã xác thực thông điệp, không phải chỗ cất mật khẩu).

Muối (salt) làm hai user cùng mật khẩu ra hash khác nhau. Thư viện passlib gắn
salt vào chuỗi lưu. Không tự viết schema salt.

\subsection{JWT, phiên và nơi cất token}

JWT cho phép máy chủ stateless: không bảng session trên RAM~\cite{rfc7519}.
Nhược: thu hồi sớm khó hơn cookie phiên phía server (đồ án chưa denylist
\texttt{jti}). Access ngắn (15 phút) giảm cửa sổ lộ; refresh dài (7 ngày) giữ
đăng nhập. Hai loại token phân biệt bằng claim \texttt{type} để không lấy
refresh gọi API nghiệp vụ.

\texttt{localStorage} dễ đọc bằng JavaScript: nếu có XSS, token bị lấy. Cookie
\texttt{HttpOnly} chịu XSS tốt hơn nhưng phức tạp CSRF. Đồ án chọn
\texttt{localStorage} cho đồ án ngành, ghi rõ hướng chuyển cookie khi triển khai
công cộng.

\subsection{XSS}

XSS (Cross-Site Scripting) chèn mã vào trang của nạn nhân. React mặc định escape
text trong JSX. Nguy cơ còn lại: \texttt{dangerouslySetInnerHTML}, Markdown không
lọc, URL javascript:. Đồ án không render HTML thô từ câu trả lời mô hình; tin
nhắn là text. Highlight tìm kiếm tách token rồi bọc phần tử, không ghép HTML
chuỗi.

\subsection{CSRF}

CSRF (Cross-Site Request Forgery) lừa trình duyệt đã đăng nhập gửi request sang
site đích. Cookie tự đính; header \texttt{Authorization} do JS gắn thì form giả
từ origin khác không tự có Bearer. Mô hình token đồ án giảm CSRF so với cookie
phiên cổ điển. Khi chuyển cookie, phải thêm token CSRF hoặc SameSite.

\subsection{Giới hạn tần suất}

Rate limiting chặn một client gọi quá dày. Với RAG, động cơ không chỉ DoS máy
chủ mà còn vét hạn mức API trả phí. slowapi đếm theo user (hoặc IP). Giới hạn
quá chặt làm demo hỏng; quá lỏng thì một vòng lặp client đốt hết hạn mức ngày.

\section{Tệp văn bản, hàng đợi nền và kiến trúc 12 yếu tố}

\subsection{PDF}

PDF mô tả trang in: chữ, font, toạ độ, đôi khi chỉ ảnh quét. Thư viện trích chữ
đọc lớp text. File scan không có lớp chữ thì trích ra rỗng --- đồ án từ chối lúc
tải, không đợi soi hợp đồng mới báo. Không OCR trong phạm vi đồ án vì tốn GPU
và sai trên biểu mẫu.

\subsection{DOCX}

DOCX là gói XML trong ZIP (Office Open XML). Trích chữ qua thư viện đọc
\texttt{document.xml}. Định dạng lạ (doc nhị phân cũ, bảng phức tạp) có thể mất
cấu trúc; đồ án nhận DOCX phổ biến, không cam kết mọi máy gõ luật sư.

\subsection{Việc chạy nền}

Soát xét hợp đồng gọi LLM nhiều lần (từng cụm điều khoản), không phù hợp giữ
một request HTTP đến khi xong nếu người dùng đóng tab. FastAPI
\texttt{BackgroundTasks} chạy trong cùng tiến trình: đủ đồ án, mất việc nếu
container chết giữa chừng. Hàng đợi phân tán (Redis/Celery) là hướng nâng cấp,
chưa cài.

\subsection{Biến môi trường và bí mật}

Mười hai yếu tố (twelve-factor) khuyên cấu hình qua môi trường, không gắn cứng
trong mã. Đồ án: khóa API, JWT, chuỗi Postgres nằm \texttt{.env}; hành vi RAG
nằm \texttt{config.yaml}. Compose chỉ đọc \texttt{.env} gốc. Nhầm file làm backend
trong Docker không thấy khóa, hỏi đáp 503 --- đã gặp khi triển khai.

\section{Kiểm thử phần mềm và quản lý phiên bản}

\subsection{Các cấp kiểm thử}

Kiểm thử đơn vị (unit) nhắm một hàm (parser SSE, công thức hạn chót). Kiểm thử
tích hợp nhắm nhiều lớp (API + Postgres). Kiểm thử giao diện đầu cuối nhắm trình
duyệt. Đồ án: pytest tích hợp trên Postgres tạm; Vitest đơn vị parser SSE. Chưa
có bộ Playwright --- hạn chế ghi ở Chương~5.

\subsection{Giá trị kiểm thử với RAG}

Test không thể khẳng định mọi câu hỏi pháp lý đúng luật. Test khẳng định: đăng
nhập cấp đúng loại token; user A không đọc hội thoại B; FTS không dấu trả Điều
còn dấu; thiếu khóa không làm crash startup; schema Pydantic từ chối body sai.
Đó là hàng rào hồi quy, không thay luật sư.

\subsection{Git}

Git lưu lịch sử thay đổi dạng đồ thị commit. Nhánh, diff, blame phục vụ hoàn
tác và đối chiếu khi sửa schema. Đồ án dùng Git trong quá trình phát triển. Báo
cáo không phụ thuộc việc có remote công khai.

\section{Tóm tắt chỗ gắn lý thuyết vào sản phẩm}

Bảng~\ref{tab:ly-thuyet-gan} không lặp lại bảng công nghệ Chương~2 mà chỉ ra
\emph{khái niệm nào} quyết định \emph{hành vi nào} người dùng thấy.

\begin{table}[htp]
\centering
\caption{Ánh xạ khái niệm nền --- hành vi hệ thống.}
\label{tab:ly-thuyet-gan}
\begin{tabular}{p{4.2cm}p{9.5cm}}
\toprule
\textbf{Khái niệm} & \textbf{Hành vi trong đồ án} \\
\midrule
HTTP không trạng thái & Mỗi API tự mang JWT; server không nhớ phiên RAM. \\
SSE trên HTTP/1.1 & Ô chat hiện từng bước và từng chữ. \\
JSON + Pydantic & Lỗi 422 có cấu trúc, TypeScript bám tên trường. \\
CORS khoá origin & Dev hai cổng vẫn gọi được; không mở \texttt{*}. \\
UTF-8 / unaccent & Gõ không dấu, đọc còn dấu. \\
ACID + khóa ngoại & Xóa user không để hội thoại mồ côi (cascade đã thiết kế). \\
GIN / B-tree & Cây luật và ô tìm mở được trên \soDieu{} Điều. \\
Argon2 / JWT & Đăng nhập; mật khẩu không lưu rõ. \\
XSS escape JSX & Câu LLM không thành mã chạy trong trình duyệt. \\
Background task & Tải xong PDF, soát xét chạy tiếp khi rời trang. \\
\bottomrule
\end{tabular}
\end{table}
